\begin{center}
\bfseries \fontsize{24}{12} \selectfont Abstract
\end{center}
\vspace*{1\baselineskip}

\noindent Modern cloud-native applications are built on \textbf{microservice architectures} comprising dozens to hundreds of loosely coupled services. While this architecture provides agility and scalability, it introduces immense complexity in \textbf{incident management and root cause analysis (RCA)}. A single service degradation can cascade through a mesh of dependencies, manifesting as a visible failure far from its origin. Traditional monitoring tools relying on threshold-based alerting and manual investigation struggle to isolate root causes in such environments, often resulting in \textbf{Mean Time to Resolution (MTTR)} exceeding 30--60 minutes per incident. To address this critical gap, we developed \textbf{SYNAPSE} — an AI-powered AIOps platform for \textbf{automated Root Cause Analysis of microservice failures}. SYNAPSE integrates a ten-stage AI pipeline: real-time metrics ingestion from Prometheus, service dependency graph construction, \textbf{VAE-GNN anomaly detection} (Variational Autoencoder Graph Neural Network with GAT architecture), four-tier \textbf{severity classification} (L1--L4), \textbf{DECI causal discovery} constrained by a Retrieval-Augmented Generation (RAG) prior, \textbf{LLM-powered explanation} via GPT-4o-mini, automated \textbf{recovery action selection} with simulated kubectl execution, \textbf{continual learning} via Elastic Weight Consolidation (EWC) and Model-Agnostic Meta-Learning (MAML), and persistent result storage. Evaluation on synthetic fault injection scenarios — including latency spikes, error surges, CPU throttling, memory leaks, and network partitions — across a 10-service Kubernetes cluster demonstrates a \textbf{root cause identification accuracy of 87\%}, an \textbf{average end-to-end pipeline latency of 8.3 seconds}, and a \textbf{GNN anomaly detection F1-score of 0.91}. The platform's React 18 dashboard provides live WebSocket anomaly streaming, interactive causal graph visualization, and incident management. SYNAPSE's modular, containerized architecture supports deployment on local workstations or Kubernetes clusters and is designed to \textbf{learn continuously} without forgetting previously observed fault patterns.

\vspace{\baselineskip}
\noindent \textbf{Keywords: } \textit{AIOps, Root Cause Analysis, Microservices, Graph Neural Network, VAE-GNN, Causal Discovery, DECI, LLM Reasoning, Continual Learning, EWC, MAML, Kubernetes, Anomaly Detection, Incident Management}
